\chapter{光滑流形}
\label{chap:smooth-manifolds}

正则曲线、正则曲面和正则水平集都有一个共同特征：在每一点附近，可以用
若干实参数作坐标，但整个对象通常不能由一张坐标图覆盖。球面上的经纬坐标
在极点失效，并不说明球面在那里不光滑；需要更换坐标，而不是删除极点。
流形的定义把局部坐标及坐标之间的变换纳入对象本身。

本章讨论有限维、Hausdorff、第二可数且不带边界的光滑流形。带边界流形
将在第\ref{chap:stokes-manifolds}章随 Stokes 定理引入。本章最后证明紧
光滑流形可以嵌入某个有限维欧氏空间；一般的 Whitney 嵌入定理则需要这里
尚未建立的构造。

\section{拓扑流形和坐标图}

\begin{definition}
设 \(n\in\N\)。一个 \(n\) 维拓扑流形是一个 Hausdorff、第二可数的
拓扑空间 \(M\)，并且每个 \(p\in M\) 都有开邻域 \(U\) 和到某个欧氏
开集 \(V\subset\R^n\) 的同胚
\[
 x:U\longrightarrow V.
\]
二元组 \((U,x)\) 称为一张\emph{坐标图}，\(x^1,\ldots,x^n\) 称为
这张图中的坐标函数。
\end{definition}

局部欧氏条件提供计算所需的坐标。Hausdorff 条件保证紧集闭且极限唯一；
第二可数性则使坐标图能够选成可数族，并在后面导出可数、局部有限的单位
分解。这两个条件都不由局部欧氏性自动推出。

删去其中任一条件会出现很具体的反例。取两份实直线，并把两份中所有非零点
按相同坐标粘合，只保留两个不同的原点，所得“有两个原点的直线”局部同胚于
$\R$ 且第二可数，却不能把两个原点放进不相交邻域，因此不是 Hausdorff 空间。
另一方面，不可数多份实直线的不交并是 Hausdorff 且局部欧氏的，却不第二
可数：每个分支都是非空开集，而一组可数基不可能为不可数个两两不交的开分支
各提供一个包含于其中的基元素。

欧氏开集本身是流形，恒等映射就是坐标图。更能说明多张坐标图必要性的例子是
球面
\[
 S^n=\{(u,t)\in\R^n\times\R:\norm{u}^2+t^2=1\}.
\]
从北极和南极作立体投影，得到覆盖 \(S^n\) 的两张图。两图重叠处的坐标
变换为
\[
 u\longmapsto \frac{u}{\norm{u}^2},\qquad u\ne0.
\]
它与逆映射在各自定义域上连续，因而球面是 \(n\) 维拓扑流形。

当 \(n\ge1\) 时，\(S^n\) 不可能由一张坐标图覆盖。否则它将同胚于
\(\R^n\) 的一个非空开集；前者紧，而后者不紧，矛盾。坐标图是局部描述，要求
整个流形只有一张图会排除最基本的紧流形。

积空间给出另一个基本构造。若 \((U,x)\)、\((V,y)\) 分别是 \(M^m\)、
\(N^n\) 的坐标图，则
\[
 (x,y):U\times V\longrightarrow x(U)\times y(V)\subset\R^{m+n}
\]
是积空间的坐标图。Hausdorff 性和第二可数性在有限乘积下保持，所以
\(M\times N\) 是 \(m+n\) 维拓扑流形。

定义中预先固定了维数。一个非空连通空间不可能既是 \(m\) 维流形又是
\(n\) 维流形而 \(m\ne n\)。这一事实最终依赖 Brouwer 不变域定理；
本书不在此处证明它。

\section{图册、相容性和光滑结构}

同一开集可以使用许多坐标。要在不同坐标中谈微分，必须限制坐标变换的
正则性。设 \((U,x)\)、\((W,y)\) 是两张拓扑坐标图。若
\(U\cap W\ne\varnothing\)，则
\[
 y\circ x^{-1}:x(U\cap W)\longrightarrow y(U\cap W)
\]
是欧氏开集之间的同胚。若它和逆变换都属于 \(C^\infty\)，便称这两张图
\emph{光滑相容}；交集为空时约定它们相容。

覆盖 \(M\) 的一族两两光滑相容坐标图称为光滑图册。一个光滑图册与所有
同它相容的坐标图合在一起，得到唯一的极大光滑图册；极大光滑图册称为
\(M\) 上的光滑结构。这里还须验证新增的图彼此相容：在任意两张新图
的交点附近，取原图册中覆盖该点的一张图 $z$，将新图间的变换写成
$(y\circ z^{-1})\circ(z\circ x^{-1})$。两个因子光滑，故变换在每个交点
附近光滑；光滑性是局部性质，因而在整个交集上成立。赋予光滑结构的拓扑
流形称为光滑流形。

\begin{proposition}\label{prop:atlas-equivalence}
两个光滑图册决定同一光滑结构，当且仅当其中每张图都与另一个图册中的
每张图光滑相容。
\end{proposition}

\begin{proof}
若两图册生成同一极大图册，结论来自极大图册中任意两张图的相容性。
反过来，若两个图册交叉相容，则它们的并仍是光滑图册。若一张新图与第一
图册相容，在它与第二图册任意一张图的交点附近，取第一图册中覆盖该点的图
作为中间坐标。上述局部分解说明新图也与第二图册相容。交换两图册同理，
所以它们所有相容图组成的极大图册相同。
\end{proof}

立体投影的坐标变换不仅连续，而且在
\(\R^n\setminus\{0\}\) 上光滑，故给 \(S^n\) 一个光滑结构。若
\(M,N\) 已有光滑结构，积坐标之间的变换分别由 \(M,N\) 上的坐标变换
组成，因而给出积光滑结构。

光滑结构不是单张坐标图，也不依赖于预先把流形画在某个欧氏空间中。它记录
哪些坐标可以光滑地互换。切向量、微分形式和积分都必须在相容坐标变换下
表示同一个对象。

\section{光滑映射和微分同胚}

设 \(M^m,N^n\) 为光滑流形，\(F:M\to N\) 为连续映射。取 \(p\in M\)
附近的坐标图 \((U,x)\) 和 \(F(p)\) 附近的坐标图 \((V,y)\)，并缩小
\(U\) 使 \(F(U)\subset V\)。映射在坐标中的表示为
\[
 y\circ F\circ x^{-1}:x(U)\longrightarrow y(V).
\]
若对每个 \(p\) 都可取坐标使这个表示为 \(C^r\) 映射，则称 \(F\) 是
\(C^r\) 映射；\(r=\infty\) 时称 \(F\) 光滑。连续性也可不预先写入
定义，而要求上述局部坐标表示存在并且为 \(C^r\)；两种定义等价。

这个定义不依赖所选坐标。若换成
\((\widetilde U,\widetilde x)\) 和 \((\widetilde V,\widetilde y)\)，
新的坐标表示在公共定义域上等于
\[
 (\widetilde y\circ y^{-1})\circ
 (y\circ F\circ x^{-1})\circ
 (x\circ\widetilde x^{-1}).
\]
两侧附加的坐标变换光滑，故中间一项的 \(C^r\) 性在换图后保持。

\begin{proposition}\label{prop:smooth-map-calculus}
光滑流形的恒等映射光滑，光滑映射的复合仍光滑。若 \(F:M\to N\) 是双射，
并且 \(F,F^{-1}\) 都光滑，则称 \(F\) 为微分同胚。
\end{proposition}

\begin{proof}
恒等映射在同一张坐标图中表示为欧氏空间的恒等映射。复合映射在适当坐标中
是两个欧氏光滑映射的复合，结论来自多元链式法则。最后一句是定义；应当注意，
光滑双射的逆映射未必光滑。
\end{proof}

最简单的例子是 \(F:\R\to\R\)，\(F(t)=t^3\)。它是光滑双射，但逆映射
\(s\mapsto s^{1/3}\) 在零点不可微。若 \(DF\) 处处可逆，逆函数定理才保证
局部逆光滑；再结合双射性，才可拼成全局微分同胚。

坐标表示的导数秩也与坐标无关。更换源坐标和目标坐标时，Jacobian 矩阵
分别在右侧和左侧乘以可逆矩阵，所以秩不变。若 \(F:M^m\to N^n\) 在
每一点的坐标导数秩都为 \(m\)，称 \(F\) 为\emph{浸入}；若秩都为
\(n\)，称 \(F\) 为\emph{浸没}。浸入只说明局部切方向没有被压缩为零，
并不保证整体单射，也不保证像集带着正确的子空间拓扑。

映射 \(t\mapsto(\cos t,\sin t)\) 从 \(\R\) 到 \(S^1\) 是处处秩为一的浸入，
但具有周期，因而不单射。即使另加单射仍然不够，下面的例子说明还需控制像上
的拓扑。

例如
\[
 \gamma:(-\pi,\pi)\longrightarrow\R^2,\qquad
 \gamma(t)=(\sin t,\sin 2t)
\]
是单射浸入，但不是到其像的同胚。当 \(t_k\uparrow\pi\) 时，
\(\gamma(t_k)\to(0,0)=\gamma(0)\)，而 \(t_k\not\to0\)。因此，定义
光滑嵌入时还必须加入拓扑条件：一个光滑映射若是浸入，又是到其像（赋予
子空间拓扑）的同胚，就称为\emph{光滑嵌入}。

\section{子流形和正则值}

\begin{definition}
设 \(M^n\) 为光滑流形，\(S\subset M\)。若存在固定的
\(0\le k\le n\)，使对每个 \(p\in S\) 都有 \(M\) 的坐标图
\((U,x)\)，满足 \(x(p)=0\) 以及
\[
 x(U\cap S)=x(U)\cap\bigl(\R^k\times\{0\}\bigr),
\]
则称 \(S\) 为 \(M\) 的 \(k\) 维嵌入子流形。这里 \(S\) 取 \(M\)
的子空间拓扑。
\end{definition}

这样的坐标称为切片坐标。坐标限制
\[
 x_S:U\cap S\longrightarrow
 \{u\in\R^k:(u,0)\in x(U)\},\qquad
 x_S=(x^1,\ldots,x^k),
\]
构成 \(S\) 的光滑图册。两张切片坐标之间的变换是原坐标变换在切片上的
限制，因而光滑。这给 \(S\) 唯一的光滑结构，使所有切片坐标光滑，并使
包含映射 \(\iota:S\hookrightarrow M\) 成为光滑嵌入。

图像是最直接的例子。若 \(U\subset\R^k\) 开且
\(f:U\to\R^{n-k}\) 光滑，则
\[
 \Gamma_f=\{(u,f(u)):u\in U\}
\]
是 \(k\) 维嵌入子流形；坐标变换
\[
 (u,v)\longmapsto (u,v-f(u))
\]
把它化为 \(v=0\) 的切片。

\begin{theorem}[正则值定理]\label{thm:regular-value-submanifold}
设 \(F:M^n\to N^m\) 光滑，\(m\le n\)，\(c\in N\)。若对每个
\(p\in F^{-1}(c)\)，\(F\) 在 \(p\) 的坐标导数都是满射，则
\(F^{-1}(c)\) 是 \(M\) 的 \(n-m\) 维嵌入子流形。
\end{theorem}

\begin{proof}
固定 \(p\in F^{-1}(c)\)，在 \(p,c\) 附近分别取坐标 \(x,y\)，并平移
坐标使 \(x(p)=0\)、\(y(c)=0\)。坐标表示
\[
 \widetilde F=y\circ F\circ x^{-1}
\]
在 \(0\) 处的导数满射。第\ref{chap:local-inversion}章的正则水平集定理
给出源空间中的局部坐标变换 \(G\)，使
\[
 \widetilde F\circ G^{-1}(u,v)=v,\qquad
 (u,v)\in\R^{n-m}\times\R^m.
\]
于是 \(F^{-1}(c)\) 在坐标 \(G\circ x\) 中恰是 \(v=0\) 的切片。
每一点都有这样的切片坐标，故它满足嵌入子流形的定义。
\end{proof}

当 $n\ge1$ 时，单位球面 $S^{n-1}\subset\R^n$ 是函数
\(F(x)=\norm{x}^2\) 的水平集 \(F^{-1}(1)\)。在球面上
\[
 DF(x)h=2x\cdot h
\]
是满射，所以正则值定理重新得到 \(S^{n-1}\) 的光滑结构。相反，
\(F(x,y)=x^2-y^2\) 的零水平集在原点交叉，原点不是正则点；这一水平集
不能在原点附近化成一条直线。正则值条件是水平集成为子流形的充分条件，
却不是必要条件：\(F(x)=x^3\) 的零水平集是单点流形，但 \(0\) 不是
正则值。

\section{单位分解与嵌入}

设 \(\{U_\alpha\}_{\alpha\in A}\) 是 \(M\) 的开覆盖。光滑函数族
\(\{\rho_i\}_{i\in I}\) 称为从属于这个覆盖的光滑单位分解，如果
\[
 0\le\rho_i\le1,\qquad
 \{\operatorname{supp}\rho_i\}_{i\in I}\ \text{局部有限},\qquad
 \sum_{i\in I}\rho_i=1,
\]
并且每个 \(i\) 都有 \(\alpha(i)\in A\) 使
\(\operatorname{supp}\rho_i\subset U_{\alpha(i)}\)。局部有限性保证
每一点附近只有有限项非零，因此这个和是光滑函数，并且可以逐项微分。

先记录欧氏截断函数在流形上的版本。

\begin{lemma}[光滑截断]\label{lem:manifold-smooth-cutoff}
若 \(K\subset M\) 紧，\(U\subset M\) 开且 \(K\subset U\)，则存在
\(\chi\in C_c^\infty(M)\) 满足
\[
 0\le\chi\le1,\qquad
 \chi=1\ \text{于 \(K\) 的一个邻域},\qquad
 \operatorname{supp}\chi\subset U.
\]
\end{lemma}

\begin{proof}
先设 \(K\) 包含在一张坐标图 \((V,x)\) 中。缩小坐标邻域，可取开集
\(W\) 使
\[
 K\subset W,\qquad \overline W\subset V\cap U,
\]
其中 \(\overline W\) 紧。欧氏空间中的光滑截断函数给出
\(\eta\in C_c^\infty(x(V\cap U))\)，使 \(0\le\eta\le1\)，并在
\(x(\overline W)\) 的邻域等于 \(1\)。把 \(\eta\circ x\) 定义在
\(V\cap U\) 上，并在 \(M\setminus(V\cap U)\) 上置零，便得到所需函数。
由于其支集是 \(V\cap U\) 的紧子集，补零处附近函数恒为零，因而补零后
仍光滑。

一般情形下，先取有限多个坐标图 \(V_j\) 及开集 \(W_j\)，使
\(W_1,\ldots,W_r\) 覆盖 \(K\)，并且
\(\overline W_j\Subset V_j\cap U\)。对每个紧集 \(\overline W_j\) 作上述构造，
所得截断函数记为 \(\chi_1,\ldots,\chi_r\)。令
\[
 \chi=1-\prod_{j=1}^r(1-\chi_j).
\]
它取值于 \([0,1]\)，在 \(K\) 的一个邻域等于 \(1\)，而其支集包含在
有限个 \(\operatorname{supp}\chi_j\) 的并中，故为 \(U\) 的紧子集。
\end{proof}

截断引理还能把局部函数延拓成全局函数。若 \(f\in C^\infty(U)\)，而
\(K\subset U\) 紧，取上述 \(\chi\) 并在 \(U\) 上令
\(\widetilde f=\chi f\)，在 \(M\setminus U\) 上令 \(\widetilde f=0\)。
由于 \(\operatorname{supp}\chi\) 紧含于 \(U\)，补零不会在 \(\partial U\)
产生光滑性问题；所得 \(\widetilde f\in C_c^\infty(M)\) 在 \(K\) 的一个
邻域与 \(f\) 相同。这种操作使坐标中的局部构造可以先在较小紧集上保留，再
放到整个流形上。

\begin{theorem}[光滑单位分解]\label{thm:smooth-partition-of-unity}
光滑流形上的每个开覆盖都有从属于它的可数光滑单位分解。
\end{theorem}

\begin{proof}
先构造紧穷竭。由第二可数性，可以选取可数个闭包紧的坐标邻域
\(B_1,B_2,\ldots\) 覆盖 \(M\)。局部紧性还说明：每个紧集都有闭包紧的
开邻域。递归地选择紧集 \(K_j\)，使
\[
 \overline{B_1}\subset\operatorname{int}K_1,\qquad
 K_j\cup\overline{B_{j+1}}\subset\operatorname{int}K_{j+1}.
\]
于是
\[
 K_1\subset\operatorname{int}K_2\subset K_2
 \subset\operatorname{int}K_3\subset\cdots,\qquad
 M=\bigcup_{j\ge1}K_j.
\]
约定 \(K_j=\varnothing\) 当 \(j\le0\)。

紧层 \(K_j\setminus\operatorname{int}K_{j-1}\) 包含在开集
\(\operatorname{int}K_{j+1}\setminus K_{j-2}\) 中。对其中每一点选择
开坐标邻域 \(W_{j\ell}\) 及稍大的坐标邻域 \(V_{j\ell}\)，使
\[
 \overline{W_{j\ell}}\Subset V_{j\ell}
 \subset U_{\alpha(j,\ell)}\cap
 \bigl(\operatorname{int}K_{j+1}\setminus K_{j-2}\bigr).
\]
由紧性，每层只需有限多个 \(W_{j\ell}\)。这些 \(W_{j\ell}\) 覆盖
\(M\)，而族 \(\{V_{j\ell}\}\) 局部有限。事实上，若一个邻域包含在
\(\operatorname{int}K_r\) 中，它不可能遇到指标 \(j\ge r+2\) 的
\(V_{j\ell}\)，其余指标只有有限多个。

由光滑截断引理，为每对指标取 \(0\le\psi_{j\ell}\le1\)，使它在
\(\overline{W_{j\ell}}\) 附近等于 \(1\)，且
\(\operatorname{supp}\psi_{j\ell}\subset V_{j\ell}\)。局部有限和
\[
 s=\sum_{j,\ell}\psi_{j\ell}
\]
光滑并处处严格为正。因此
\[
 \rho_{j\ell}=\frac{\psi_{j\ell}}{s}
\]
构成所求单位分解。除以处处正的 \(s\) 不改变支集，也不破坏局部有限性。
\end{proof}

单位分解允许把局部对象相加，但这种粘合要求目标具有线性或凸结构。例如局部
函数和局部内积可以加权平均；局部微分同胚加权以后一般不再可逆。

\begin{theorem}[紧流形的有限维嵌入]\label{thm:compact-manifold-elementary-embedding}
每个紧光滑 \(n\) 维流形都可光滑嵌入某个有限维欧氏空间。
\end{theorem}

\begin{proof}
取有限坐标覆盖 \(\{(U_i,x_i)\}_{i=1}^m\)，并对这个覆盖使用单位分解
定理。所得分解局部有限，而紧空间上的局部有限非空集族必为有限族：取每点
只遇有限个支集的邻域，再由紧性取有限子覆盖即可。删去恒为零的成员，并把
隶属于同一 \(U_i\) 的函数相加，可设单位分解恰写成
\(\{\rho_i\}_{i=1}^m\)，并且
\(\operatorname{supp}\rho_i\subset U_i\)。函数
\(\rho_i x_i:M\to\R^n\) 先在 \(U_i\) 上按乘积定义，再在 \(U_i\)
外取零。由于 \(M\) 紧，\(\operatorname{supp}\rho_i\) 是 \(U_i\)
中的紧集，故补零后光滑。定义
\[
 \Phi(p)=
 \bigl(\rho_1(p),\rho_1(p)x_1(p),\ldots,
       \rho_m(p),\rho_m(p)x_m(p)\bigr)
 \in\R^{m(n+1)}.
\]

若 \(\Phi(p)=\Phi(q)\)，选 \(i\) 使 \(\rho_i(p)>0\)。对应的标量分量
说明 \(\rho_i(q)=\rho_i(p)>0\)，所以 \(p,q\in U_i\)；向量分量再给出
\(\rho_i(p)x_i(p)=\rho_i(q)x_i(q)\)，从而 \(x_i(p)=x_i(q)\)，故
\(p=q\)。因此 \(\Phi\) 单射。

再证它是浸入。在 \(p\) 处选 \(i\) 使 \(\rho_i(p)>0\)，并用 \(x_i\)
作为源坐标。若坐标方向 \(v\in\R^n\) 被
\(D(\Phi\circ x_i^{-1})(x_i(p))\) 映为零，则第 \(i\) 个标量分量给出
\(D\rho_i(p)v=0\)，而相应向量分量给出
\[
 D(\rho_i x_i)(p)v
 =\rho_i(p)v+x_i(p)D\rho_i(p)v
 =\rho_i(p)v=0.
\]
由于 \(\rho_i(p)>0\)，必有 \(v=0\)，所以坐标导数的秩为 \(n\)。
最后，紧空间到 Hausdorff 空间的连续单射是到其像的同胚，故 \(\Phi\)
是光滑嵌入。
\end{proof}

这个证明给出的欧氏空间维数取决于所选有限图册，并不声称最佳界。一般非紧
流形也能嵌入有限维欧氏空间，但完整的 Whitney 嵌入定理不能由上述有限覆盖
论证得到，故不在本章证明范围内。

\section*{习题}
\addcontentsline{toc}{section}{习题}

\subsection*{A组　坐标与光滑结构}
\begin{enumerate}[label=\textbf{\thechapter.\arabic*.}]
 \exerciseitem{manifold-sphere-stereographic}
 写出 \(S^n\) 两张立体投影的逆映射，并直接验证重叠处坐标变换光滑。
 \exerciseitem{manifold-product-atlas}
 证明两个光滑流形的积图构成光滑图册，并验证三重乘积的两种括号方式由自然
 微分同胚联系。
 \exerciseitem{manifold-atlas-common-refinement}
 若一个图册的坐标域同时加细两个给定图册的坐标域覆盖，且其中各图与两个
 给定图册都光滑相容，就称它为共同光滑加细。证明两个光滑图册等价，当且
 仅当它们有一个共同光滑加细。
 \exerciseitem{manifold-smoothness-all-charts}
 证明光滑映射若在覆盖源、目标的两组图中具有光滑坐标表示，则它在相应极大
 图册的任意坐标中都光滑。
 \exerciseitem{manifold-rank-coordinate-invariance}
 用链式法则证明映射坐标导数的秩与坐标选择无关。
\end{enumerate}

\subsection*{B组　子流形与正则值}
\begin{enumerate}[label=\textbf{\thechapter.\arabic*.},resume]
 \exerciseitem{manifold-graph-submanifold}
 完整写出把光滑函数图像拉直的坐标变换，并证明图像到定义域的投影是
 微分同胚。
 \exerciseitem{manifold-sphere-regular-value}
 设 $n\ge1$。用正则值定理证明 \(S^{n-1}\subset\R^n\) 是嵌入子流形，并在一点写出
 局部图像坐标。
 \exerciseitem{manifold-special-linear-submanifold}
 设 \(n\ge1\)。证明
 \[
 SL(n,\R)=\{A\in\R^{n\times n}:\det A=1\}
 \]
 是 \(\R^{n^2}\) 中的 \(n^2-1\) 维嵌入子流形。
 \exerciseitem{manifold-singular-level-examples}
 分别分析 \(x^2-y^2=0\) 与 \(x^3=0\) 的零水平集，说明正则值条件失效
 时水平集可能不光滑，也可能仍是子流形。
 \exerciseitem{manifold-injective-immersion-not-embedding}
 验证
 \[
 \gamma(t)=(\sin t,\sin 2t),\qquad -\pi<t<\pi,
 \]
 是单射浸入但不是嵌入。
\end{enumerate}

\subsection*{C组　单位分解与有限维嵌入}
\begin{enumerate}[label=\textbf{\thechapter.\arabic*.},resume]
 \exerciseitem{manifold-cutoff-finite-combination}
 补全光滑截断引理中有限个局部截断合成全局截断的支集与光滑性论证。
 \exerciseitem{manifold-locally-finite-smooth-sum}
 证明局部有限光滑函数族的逐点和光滑，并给出“逐点只有有限项非零”但不局部
 有限时和函数不连续的例子。
 \exerciseitem{manifold-partition-glue-functions}
 设 \(f_i\in C^\infty(U_i)\) 在重叠处不必相等。证明
 \(f=\sum_i\rho_i f_i\) 良定义且光滑，并说明它通常不在任何 \(U_i\)
 上等于 \(f_i\)。
 \optionalexerciseitem{manifold-compact-embedding-reconstruction}
 不引用本章最后一个定理，使用有限坐标覆盖和单位分解重新构造紧流形到
 \(\R^{m(n+1)}\) 的嵌入，并逐项验证单射、浸入和拓扑嵌入。
\end{enumerate}
